\Needspace{7\baselineskip}
\subsection{The positive model on arbitrary modules}\label{sec:positive-laws}
The usual divided powers do satisfy the mixed identities, including for
modules that are neither flat nor finitely generated. We spell out this
positive model to distinguish it from the abstract counterexamples.

Roby's representing theorem gives
\[
 \operatorname{Pol}^d_k(X,Y)=\Hom_k(\Gamma^d_kX,Y)
\]
for arbitrary modules over a commutative ring. His scalar-extension and
exponential theorems give, respectively,
\[
 R\otimes_k\Gamma_k^dX\simeq\Gamma_R^d(R\otimes_kX),\qquad
 \Gamma^d(X\oplus Y)\simeq
 \bigoplus_{a+b=d}\Gamma^aX\otimes_k\Gamma^bY.
\]
These are \cite[Theorem~IV.1, pp.~266--267; Theorems~III.3--III.4,
pp.~262--263]{Roby}. In particular, invariant tensors must not replace
the universal divided powers on an arbitrary nonflat module.

\begin{proposition}\label{prop:ordinary-polynomial-category}
The modules over $k$, with morphisms
$\bigoplus_{i\ge0}\Hom_k(\Gamma^iX,Y)$, form the category of
polynomial laws with a globally finite degree bound. Composition is
composition after every scalar extension. If
$f=\sum_{i=0}^m f_i$ and $g_q:\Gamma^qY\to Z$, its degree-$n$ part is
\begin{equation}\label{eq:ordinary-mixed}
 \sum_{\substack{\sum_i\nu_i=q\\\sum_i i\nu_i=n}}
 g_q\mu_\nu
 \left(\bigotimes_i\Gamma^{\nu_i}(f_i)\right)
 \left(\bigotimes_i\alpha_{\nu_i,i}\right)
 \Delta_{(i\nu_i)_i}.
\end{equation}
Here $\alpha_{a,b}:\Gamma^{ab}X\to\Gamma^a\Gamma^bX$
represents $x\mapsto\gamma_a(\gamma_b(x))$, $\Delta$ is the
specified component of the divided-power diagonal, and $\mu_\nu$
is multiplication into $\Gamma^qY$. The homogeneous degree-one subcategory is
$\mathsf{Mod}_k$, and finite products are direct sums.
\end{proposition}
\begin{proof}
For a homogeneous pair, naturality of the universal law gives
$g_q\Gamma^q(f_p)\alpha_{q,p}$. For a sum of homogeneous laws,
apply the divided-power identity
\[
 \gamma_q(y_0+\cdots+y_m)=
 \sum_{\nu_0+\cdots+\nu_m=q}
 \gamma_{\nu_0}(y_0)\cdots\gamma_{\nu_m}(y_m)
\]
to $y_i=f_i(\gamma_i(x))$. Its term indexed by $\nu$ has input
degree $\sum i\nu_i$. Representing that term gives exactly
\eqref{eq:ordinary-mixed}. The sum is finite; the constant term
causes no exception, since $\alpha_{q,0}(1)=\gamma_q(1)$.
Composition is associative and unital as an operation on natural
families of functions. The representing bijection transfers these
identities to the displayed formula. A pair of bounded laws has a
common finite bound and a unique law to the direct sum. Degree-one
laws are precisely linear maps, again by representation. The
exponential decomposition represents a law of multidegree
$(d_1,\ldots,d_s)$ on $X_1\oplus\cdots\oplus X_s$ by
$\bigotimes_j\Gamma^{d_j}X_j$. Thus the representing tensor
products are the usual tensor products of modules, with their
ordinary associativity, permutation and unit maps. Uniqueness of
all degree decompositions follows by substituting $tx$ over $R[t]$
and comparing powers of the indeterminate $t$.
\end{proof}

\begin{proposition}\label{prop:gamma-hom-law}
For arbitrary $k$-modules $X,Y$ and $p\ge0$, the operation
\[
 f\longmapsto\Gamma^p(f):
 \Hom_k(X,Y)\rightsquigarrow
 \Hom_k(\Gamma^pX,\Gamma^pY)
\]
is a homogeneous polynomial law of degree $p$. It preserves identities
and composition as identities of laws and is natural in both modules.
No base-change isomorphism for $\Hom$ is assumed.
\end{proposition}
\begin{proof}
For finitely many maps $f_1,\ldots,f_s$, expand
$\Gamma^p(\sum_iT_if_i)$ over $k[T_1,\ldots,T_s]$.
Its coefficient of $T^\nu$, for $\sum\nu_i=p$, is
\[
 c_\nu=\mu_\nu\left(\bigotimes_i\Gamma^{\nu_i}(f_i)\right)
          \Delta_\nu
 \in\Hom_k(\Gamma^pX,\Gamma^pY).
\]
For $\sum_i r_i\otimes f_i\in R\otimes_k\Hom_k(X,Y)$,
define the value to be $\sum_\nu r^\nu\otimes c_\nu$.
We must check that it depends only on this tensor, even when the
canonical map from the tensor product of a Hom module to a Hom module
after base change is not an isomorphism.

Adjoining a zero term, splitting either a map or a coefficient as
a sum, and balancing a $k$-scalar between its coefficient and the map generate all changes
of tensor presentation. The divided-power sum and scalar identities
show that the indicated coefficient operators respect these changes.
For the splitting of a coefficient on a repeated map, the needed
identity is $\mu_{a,b}\Delta_{a,b}=\binom{a+b}{a}\id$ on
$\Gamma^{a+b}$; it follows by expanding $\gamma_{a+b}((s+t)x)$
over $k[s,t]$ and comparing coefficients.
One can check each identity with independent indeterminates: a finite
polynomial of $k$-linear operators is zero after extension to a free
polynomial algebra only when each coefficient operator is zero.
Indeed, evaluate on any input and compare its coefficients in the
free module of monomials. Thus the coefficient equalities hold in
the ordinary Hom module itself and can be specialized to every $R$.
This proves independence of presentation and naturality in $R$.
Every monomial has degree $p$, so the law is homogeneous.
The ordinary functorial identities for $\Gamma^p$ over the free
polynomial algebra give identity and composition compatibility by
the same coefficient argument. For $p=0$ the resulting law is the
constant identity of $k$.
\end{proof}

\begin{proposition}\label{prop:three-polynomial-categories}
Over a nonzero commutative ring $k$, the category of bounded polynomial
laws, the fixed-degree category with morphisms
$\Gamma^d\Hom_k(V,W)$ on finite projective modules, and its category
of linear representations must be distinguished. The latter
representations are exactly functors whose action on linear Hom
modules is a degree-$d$ polynomial law compatible with composition
and identities. The bounded polynomial-law category is equivalent
to neither of the other two categories.
\end{proposition}
\begin{proof}
For $k=\F_2$, $V=k$, $W=k^2$ and $d=2$, the module
$\Hom_k(\Gamma^2V,W)$ has dimension two, whereas
$\Gamma^2\Hom_k(V,W)$ has dimension three. Thus even the proposed
Hom modules differ. Homogeneous polynomial laws of degree $d>1$
compose in degree $d^2$, and have no degree-one identities within
that single piece.

There is also a categorical obstruction valid for every nonzero
$k$. The bounded-law category has terminal object zero and no
initial object: every object has at least the two distinct constant
laws $0,1$ to $k$. For $d>0$, the fixed-degree category has a zero
object. For $d=0$, each of its Hom sets is $k$, so it has no
terminal object. Its representation category is abelian and has a
zero object for every $d$. These properties are invariant under
equivalence, proving the non-equivalence assertions.

Finally Roby's representation applied to the module
$\Hom_k(V,W)$ identifies a degree-$d$ law with a linear map
\[
 \Gamma^d\Hom_k(V,W)\longrightarrow\Hom_k(FV,FW).
\]
The composition and identity equations represent exactly the
functor equations in the fixed-degree category. The same argument
identifies the natural transformations. This is the standard strict
polynomial-functor category of \cite[Section~2, pp.~3--5]{Krause}.
\end{proof}
